\Needspace{18\baselineskip}
\section{Search and state isomorphism classes}
\label{sec:state-search}

The focus recovery constructions define search on actual states. Search is
deterministic up to focus-preserving isomorphism, final states have no search
step, and no search step ends in an initial state. A concrete initial value
state makes a search step to an actual final state.

We then prove source and target renaming invariance and descent to literal-label
and recursive state isomorphism classes, where successor classes are unique.
Section~\ref{sec:contextual-search} now proves equivalence of erasure-based
search with the paper's contextual local rewrites. Substitution and compute
transitions for the full machine remain to be constructed.

\Needspace{16\baselineskip}
\subsection{Search between actual states}
\begin{definition}[Erasure-based state search]\label{term:state-search}
For a fixed operation-kind function, write $s\to_{\mathrm{er}}t$ if there is a
typed coordinate $q$ on the underlying graph of $s$ such that the extracted
coordinate of $s$ makes one of the five search steps to $q$, and the realization
of $q$ is focus-preservingly isomorphic to $t$. Thus the target may use arbitrary
vertex and edge names. This is a relation on actual state records, rather than
merely a construction of two endpoints from a separately supplied graph.
\end{definition}
\lean{Graph/StateSearch}{86}{90}{Linear.State.Search}

\begin{lemma}[Existence and target renaming]
Every typed coordinate successor of an extracted state coordinate gives an
actual state successor. Replacing a successor by a focus-preservingly
isomorphic state preserves the step.
\end{lemma}
\begin{proof}
For existence, realize the successor coordinate and use the identity focussed
isomorphism as witness. For target renaming, retain the same coordinate step
and compose its target isomorphism with the supplied isomorphism.
\end{proof}
\lean{Graph/StateSearch}{93}{98}{Linear.State.search_realize}
\lean{Graph/StateSearch}{101}{104}{Linear.State.search_target_iso}

\begin{theorem}[Search determinism up to focussed isomorphism]
Two erasure-based search successors of the same state are focus-preservingly
isomorphic.
\end{theorem}
\begin{proof}
Both coordinate steps use the same erased graph and the same extracted input.
Coordinate search determinism identifies their output pairs; subtype
extensionality identifies the typed positions. Their canonical realizations
therefore coincide. Compose the inverse of the first target isomorphism with
the second target isomorphism to identify the actual successors.
\end{proof}
\lean{Graph/StateSearch}{107}{113}{Linear.State.search_deterministic}

\begin{lemma}[Explicit-state search boundary laws]
A final state has no erasure-based search step, and no such step ends in an
initial state.
\end{lemma}
\begin{proof}
If the source is final, its extracted coordinate has backtrack mode at an
input. Coordinate search cannot start there. If the target is initial, pull
initiality back through the target's focussed isomorphism and the realization
boundary equivalence. The successor coordinate would be search at an input,
which the coordinate search boundary theorem excludes.
\end{proof}
\lean{Graph/StateSearch}{116}{122}{Linear.State.final_no_search}
\lean{Graph/StateSearch}{125}{133}{Linear.State.no_search_initial}
These are search-only laws. Neither all-transition finality nor the full
robustness theorem follows until the remaining transition classes and specimen
obligations are supplied.

\Needspace{20\baselineskip}
\subsection{Concrete execution and interior recovery}
Erase the focus of the previously constructed initial value state. Its value
edge survives and has an empty argument/thunk view. Its source contracts from
the old focus target to the joined port, which is exactly the extracted state
coordinate's vertex. This proves the view's source certificate; the other
certificates follow from the nullary value signature.
\lean{Graph/ErasureExamples}{9}{22}{Linear.Examples.erasedValueOperation}
Apply the passive nullary rule, advance the typed coordinate, and realize it.
The existence lemma gives an actual state search step. The result is final
because its mode is backtrack and its vertex retains the input status of the
initial coordinate. This proves that the new state relation has a concrete
execution, rather than only universal implications about hypothetical steps.
\lean{Graph/ErasureExamples}{26}{37}{Linear.Examples.value_state_search}
For an interior test, relabel the old focus of the initial value graph as a
variable, retaining the value edge and all incidences. Both labels exclude
focuses and holes, so this is a focus-free, hole-free closed two-edge chain.
Insert search at its internal vertex to obtain a state whose focus source is
internal. It is not initial, directly from that source constructor. General
state recovery erases and reconstructs it with a focussed isomorphism, including
its distinguished focus and mode.
\lean{Graph/ErasureExamples}{40}{45}{Linear.Examples.closedChain}
\lean{Graph/ErasureExamples}{47}{49}{Linear.Examples.closedChain_focusFree}
\lean{Graph/ErasureExamples}{51}{53}{Linear.Examples.closedChain_holeFree}
\lean{Graph/ErasureExamples}{56}{58}{Linear.Examples.interiorState}
\lean{Graph/ErasureExamples}{60}{60}{Linear.Examples.interior_state_not_initial}
\lean{Graph/ErasureExamples}{64}{67}{Linear.Examples.interior_state_recovery}
